\section{Task 4}

\subsection{Iterative methods: overview and comparison}
This section analyzes three iterative methods for solving linear systems: Jacobi, Gauss--Seidel and SOR (Successive Over-Relaxation). Each transforms the system matrix \( A \) into a form suitable for iterative solution approximation.
\begin{itemize}
    \item \textbf{Jacobi method}: Split \( A \) into \( A = D + R \), where \( D \) is diagonal and \( R \) contains the off-diagonal entries. The iteration is:
    \[
    x^{(k+1)} = D^{-1}(b - Rx^{(k)}).
    \]  
    This method uses only the values of \( x^{(k)} \) from the preceding iteration.

    \item \textbf{Gauss--Seidel method}: Decompose \( A = L + U \), where \( L \) is lower triangular including the diagonal, and \( U \) is upper triangular. The iteration is:
    \[
    x^{(k+1)} = L^{-1}(b - Ux^{(k+1)}).
    \]  
    New values of \( x^{(k+1)} \) are used immediately in the current calculation, accelerating convergence.

    \item \textbf{SOR method}: This generalizes Gauss--Seidel by introducing a relaxation factor \( \omega \). The formula is:
    \[
    x^{(k+1)} = (1-\omega)x^{(k)} + \omega L^{-1}(b - Ux^{(k+1)}).
    \]  
    Choosing \( \omega \) appropriately is essential for efficiency.
\end{itemize}


\subsection{Results and conclusions:} 
For all calculations we set $x=0, b=rand(n,1) *10$, perform $1000$ trials and average their results.\\
\subsubsection{Parameter $k$}
We estimate \( k \), the number of iterations required to establish convergence. To avoid errors from a lack of diagonal dominance, we set \( c = n \cdot 10 \). This gives a dominant diagonal and removes errors due to matrix structure. We can then set \( n = 10 \), since this diagonal magnitude avoids size-related and convergence problems.

\begin{center}
\begin{tabular}{|l|l|l|l|l|l|}
\hline
Method       & 1          & 3          & 5          & 10         & 16         \\ \hline
Jacobi       & $4.0556e-03$ & $6.5079e-06$ & $9.9949e-09$ & $2.7341e-15$ & $1.1819e-17$ \\ \hline
Gauss-Siedel & $3.3702e-03$ & $4.4534e-07$ & $4.7716e-12$ & $2.6600e-17$ & $1.8228e-17$ \\ \hline
SOR          & $0.010987$   & $1.3029e-04$ & $1.2422e-06$ & $2.5774e-11$ & $2.3522e-17$ \\ \hline
\end{tabular}
\end{center}
Conclusions:
\begin{itemize}
    \item SOR: effect of a poorly chosen relaxation factor:\\
    Although SOR should theoretically converge fastest, our calculations do not confirm this. A fixed, poorly chosen relaxation factor was used, reducing its efficiency. A poor relaxation factor can therefore substantially reduce the method's effectiveness.
    \item Comparison of Jacobi and Gauss--Seidel:\\
    The results show that Gauss--Seidel converges faster from the first iteration. New values computed within an iteration are immediately used in subsequent calculations. Jacobi uses only preceding-iteration values, slowing convergence.
\end{itemize}



\subsubsection{Parameter $c$}
We estimate \( c \), the amount added to the diagonal in \( A + c \cdot \text{eye}(n) \). We set \( k = 100 \) to ensure convergence.
\begin{itemize}
    \item Case $c = n/2$
\begin{center}
\begin{tabular}{|l|l|l|l|}
\hline
Method       & $3$            & $10$           & $100$          \\ \hline
Jacobi       & $1.4872e-08$ & $2.0980e-03$ & $8.8618e-03$ \\ \hline
Gauss-Siedel & $5.8098e-16$ & $3.9518e-16$ & $17138e-16$  \\ \hline
SOR          & $5.9233e-16$ & $3.9828e-16$ & $17124e-16$  \\ \hline
\end{tabular}
\end{center}
    \item Case $c = 10*n$
    \begin{center}
\begin{tabular}{|l|l|l|l|}
\hline
Method       & $3$            & $10$           & $100$          \\ \hline
Jacobi       & $3.9344e-17$ & $2.5617e-17$ & $1.0956e-17$ \\ \hline
Gauss-Siedel & $4.0613e-17$ & $2.5840e-17$ & $1.0976e-17$ \\ \hline
SOR          & $3.7866e-17$ & $2.5573e-17$ & $1.0902e-17$ \\ \hline
\end{tabular}
    \end{center}
\end{itemize}
Conclusions:\\
\begin{itemize}
    \item The approximation of $c$ suggested by the mean $E[\sum_i^n X_i]$, where $X_1 \sim \mathrm{U}([0, 1])$, appears to be a good estimate. The Jacobi errors must arise from the algorithm itself, but I do not understand why they occur for [sentence unfinished in the original].
\end{itemize}

\subsection{Suggestions}
Stopping criterion:
Choose $p$ and define a stopping time at step $n$ when the following condition holds:
$$\left|\frac{x_n - x_{n-1}}{x_n}\right| \times 100 < p$$
